% ============================================================
% Table — Conceptual Distinctions Concerning Identity and Unity
% Book I — Ontology of Viability
% Part I — Identity
% Chapter 01 — Identity
% Section 04 — Identity and Unity
% ============================================================

\begin{table}[htbp]
\centering
\small
\caption{Conceptual Distinctions Concerning Identity and Unity}
\label{tab:identity-unity-conceptual-distinctions}

\renewcommand{\arraystretch}{1.35}

\begin{tabular}{
|p{0.19\textwidth}
|p{0.34\textwidth}
|p{0.37\textwidth}|
}
\hline
\textbf{Concept}
&
\textbf{Central Question}
&
\textbf{Distinction from Unity and Identity}
\\
\hline

Identity
&
What makes an existent the existent that it is?
&
Cannot be reduced to unity, numerical oneness, composition, or internal organization.
\\
\hline

Unity
&
In what ontological sense does an existent constitute one?
&
Cannot be identified with identity, simplicity, or the mere absence of plurality.
\\
\hline

Singularity
&
What is considered a single item or individual within a specified context?
&
Being singular in a context does not, by itself, establish the ontological unity of what is considered.
\\
\hline

Composition
&
How is an existent constituted through constituents and their relations?
&
May be relevant to unity but does not provide a universal account of it.
\\
\hline

Aggregate
&
What elements form a collection or grouping?
&
An aggregate does not automatically constitute one composed existent or one ontological unity.
\\
\hline

Structural Organization
&
How are elements, determinations, or relations organized?
&
May contribute to unity, but structural organization alone does not establish ontological unity.
\\
\hline

Individuation
&
How is something constituted or determined as a distinct existent?
&
Related to unity but not equivalent to the question of what makes something constitute one.
\\
\hline

Plurality
&
What differentiated determinations, aspects, or components coexist within an existent?
&
May coexist with unity, but does not automatically establish either unity or multiple identities.
\\
\hline

Multiplicity
&
How many entities or elements are present within a specified context?
&
Its numerical or contextual assessment does not, by itself, determine the number of ontological unities or identities.
\\
\hline

\end{tabular}

\end{table}